\section{n9. Towers, independent gauges, qualifications and original decision}
\paragraph{Supports.} G331, G345 and n1--n8.
For a finite tower $R\to B\to F$, over a top minimal $P$ write
$B\otimes_RK_P=\prod_Q L_Q$ and $h_Q=\dim_{L_Q}(F\otimes_BL_Q)$. Vector-space decomposition proves
\[
 d_P(F)=\sum_Q[L_Q:K_P]h_Q. \tag{21}
\]
Closed residue degrees multiply along closed-centre towers. These are the exact componentwise projection towers of (10), with zeros for absent components and full inseparable degrees. Constant ranks multiply; varying ranks are not replaced by a scalar.

Unnormalized traces on normal-domain towers with normal intermediate target compose on EVERY support. Choose one finite normal closure of all fields. Restrictions partition the embeddings of the largest field into intermediate embeddings and further embeddings; the purely inseparable multiplicities multiply. G331's full-support coefficient formula gives equality of the finite kernel sums, first on gauge differences and hence all continuous functions by density, then on $D$ by finite decreasing sums.

For general reduced injections, use n7 traces only under explicit intermediate admissibility: $H$ must have a full descended trace over $B$, and that result must have a full descended trace over $R$. On the normalization of each minimal $R$ component, its dominating source domain components have the indicated generic field extensions. Normalize those source components; this changes the generic field by degree one, and the normal-domain degree-one trace is exactly pullback on the whole fibre. More explicitly, if a minimal $B$ component $B/Q$ dominates $R/P$, the integral closure of $B/Q$ in its field $L_Q$ equals the integral closure of $R/P$ in $L_Q$: integrality of $B/Q$ over $R/P$ proves both containments by transitivity and the same monic equations. It contains the normalization of $R/P$ in $K_P$, and localization at its normal target centres partitions the closed source centres. Incomparability shows that any minimal $F$ component dominating $R/P$ passes through a unique minimal $B$ component over $P$; conversely each such dominating chain contributes to the composite. Thus the two intermediate normalization traces, restricted to these normal components, are the same full-support sums as their defining branch traces, with no appeal to equality merely at generic valuations. Normal field tower transitivity proves equality with the direct composite tuple on EVERY support of each normal $R$ centre. The equality also applies at lower-dimensional centres, irrespective of their zero ambient product. The iterated tuple is full-fibre constant by its assumed second admissibility; hence so is the direct tuple. Uniqueness of G331 descent gives $T_{F/R}H=T_{B/R}(T_{F/B}H)$. Direct admissibility alone does not imply either intermediate admissibility. This is not an unrestricted trace on incompatible tuples.

For reduced finite-flat towers determinant norms are transitive. On every minimal base field, decompose the finite reduced algebras into fields and use field-norm transitivity; reducedness makes this generic polynomial identity global on the base. The full-support kernels compose even at nonminimal supports: after splitting a support-field fibre, let a first Artinian factor have length $l$. Flatness identifies the successive maximal-ideal-layer quotients of the second free algebra with those of the first tensor its residue algebra. Multiplication therefore has the residue-algebra characteristic determinant repeated $l$ times. Each second residue length is multiplied by $l$, exactly the direct kernel length. Consequently all gauge moments and mass agree, density identifies kernels, and finite decreasing sums prove trace transitivity on $C,D$. Constant ranks $d,e$ give mass $de$. This is a proof of norm transitivity, not merely determinant basis independence.

Naive normalized towers fail when intermediate generic degrees vary. For EVERY characteristic let $A=k[[t]]$, $B=A\times A$, and $F=A\times k[[s]]$, mapping the second factor by $t=s^2$. These are reduced finite algebras; $[k((s)):k((t))]=2$ even in characteristic two, since Laurent exponent parity rules out $s\in k((t))$. In synchronized $t$ gauges the direct trace weights are $(1,2)$, total three. Direct average at values $(0,-1)$ is $-2/3$, while successive averages are $-1/2$. Unnormalized traces both equal $-2$. This is an exact all-characteristic algebraic example, not a finite sample claim.

Tracing several arbitrary product slots is false even flat. Let $R=k[x,y,z]_{(x,y,z)}$, $B=R\times R$ and $H=(-1,0)$, $V=(0,-1)$, genuine semilocal primary/unit profiles. Each source product of arity two has one zero input, hence is zero. But $\tau H=\tau V=-1$ and $\mu_R(-1,-1)$ has mass $e_R(\mathfrak m)=1$; normalized separate traces give mass $1/4$. Thus only ONE arbitrary slot, with the others pulled, has (13), (15), (20).

\paragraph{Exactly one gauge factor.}
For a base primary $\mathfrak b$, put $t(v)=v_{\mathfrak a}(\mathfrak b)$, $T(v)=v/t(v)$. Transport inputs by $u_{\mathfrak b}(Tv)=u_{\mathfrak a}(v)/t(v)$ and measures by
\[
 \mu_{\mathfrak b}=T_*[t\mu_{\mathfrak a}]. \tag{22}
\]
Every synchronized source lift has the same base denominator, so traces and maxima divide by it and all projections commute with (22). For an arbitrary separate source primary gauge $\mathfrak c$, let $w\in X_{B,\mathfrak c}$,
$\lambda(w)=w(\mathfrak aB)>0$ and $\rho_{\mathfrak c}(w)=w|_R/\lambda(w)$. The homogeneous pullback of $u$ is $\lambda(w)u(\rho_{\mathfrak c}w)$ and
\[
 \rho_{\mathfrak c,*}[\lambda\mu_{B,\mathfrak c}(\boldsymbol H_{\mathfrak c})]
 =\rho_{\mathfrak a,*}\mu_{B,\mathfrak aB}(\boldsymbol H_{\mathfrak a}). \tag{23}
\]
Indeed $w=w_{\mathfrak a}/w_{\mathfrak a}(\mathfrak c)$, so $\lambda=1/w_{\mathfrak a}(\mathfrak c)$ cancels the ONE G345 measure weight. In a trace slot the synchronized potential is $H_{\mathfrak a}(w/\lambda)=H_{\mathfrak c}(w)/\lambda$; therefore its trace is computed with $H_{\mathfrak c}/\lambda$, not with an unscaled slot. The other pulled slots satisfy the same homogeneous conversion, and (20) with (23) gives the exact independently gauged one-slot formula. There is no factor $\lambda^{n-1}$. Raw envelope competitors transport by the same order bijection, hence their suprema and both counterexample gaps transport. Successive gauge factors obey the intrinsic cocycle. Nonprimary ideals only have their positive-finite partial domains; no compact all-boundary theorem is asserted there.

\paragraph{Source correction retained.}
Literal SH17.4.7 prints the correction $\ell((x_1,\ldots,x_{d-1})M:_Mx_d)$ without quotienting the previous parameters. The finite correct term is
\[
 \ell(0:_{M/(x_1,\ldots,x_{d-1})M}x_d)
 =\ell\left(((x_1,\ldots,x_{d-1})M:_Mx_d)/(x_1,\ldots,x_{d-1})M\right).
\]
It follows by repeated complete 17.4.6 in the CURRENT quotient. The printed term can have infinite length. The exact PDF page 370, printed 356, was freshly rendered and viewed; the manuscript uses no new application of the false corollary. The grant qualification remains literal. All new source applications and their recursive-background boundaries are in the ledger and reproduced in this paper. No new Keeler/BFJ/Izumi application is inferred from old reading receipts.

\paragraph{Whole-proof versus whole-original decision.}
The author253 scoped proof is mathematically retained at its precise hypotheses; no counterexample to its stated positive results was found in this reconstruction. Its explicit $\texttt{candidate=false}$ history remains unchanged. The new manuscript adds n7's full conductor-admissible slot theorem, ALL of n4's quotient obstruction and the fully sourced determinant/going-down and tower/gauge qualifications. These are substantive additions, not a byte-identical republication. Source-body reading, duplicate matching, finite diagnostics and manuscript checks are separately reported; compilation is reserved to ROOT.

The exact P12.6 operator question is: for the independently geometric $C,D$ and bounded G345 product, which finite pullbacks, positive descents, envelope comparisons and trace-slot/projection laws survive on the FULL boundary, including normalization/conductor identifications? N1, n5--n8 answer the cone and bounded-product parts, and n2--n4,n6,n8 give universal corrected raw-envelope laws, sharp positive special hypotheses and unconditional no-go results. N7 supplies author253's remaining intended conductor-admissible one-slot identification. N4 resolves its proposed ambient extension issue NEGATIVELY. Arbitrary incompatible traces cannot descend universally; n3's normalization has branch traces $(-1,-2)$ at an identified fibre, already disproving that assertion.

Family 6.2 separates P12.3 existence/regularity from P12.6 compatibility, and 6.3 says construct when possible. We therefore DO NOT make solving all P12.3 or P12.5 a prerequisite for these finite-map theorems. Nor does one false strengthened normalization identity alone close the original. The precise residual issue is identification and compatibility of a SPECIFIED positive/regularized independently geometric envelope operator, beyond the raw competitor supremum and its conditional positivity transfer, in the remaining nonnormal/nonflat and closed-image settings. The grants select no such operator. A USC-regularized envelope, if intended, needs a justified regularization and a proof or obstruction for its pullback/descent; a finite quotient map need not transport neighborhood suprema as an open map. Our raw identities give no such theorem. The quotient counterexample is a genuine impossibility even when raw envelopes are positive at both ends, so that part is not an existence gap and is not listed as an unresolved extension theorem.

Accordingly the whole scoped mathematical proof verdict is PASS WITH THE STATED PUBLISHED-BACKGROUND AND SCOPE QUALIFICATIONS (a draft review judgment, not formal verification). The whole-original completion verdict is INCONCLUSIVE, $\texttt{candidate=false}$: no exact residual positive/regularized envelope calculus or complete negative classification of that remaining operator is proved. This is a narrower residual than author253's broad request for all envelope existence and arbitrary incompatible operator extension. All nilpotent-generic, unbounded-energy, arbitrary-model, plurifine and literal radial theories remain outside the grants; they are not silently counted as new P12.6 obligations. Smooth specialization remains exactly CLOSED affine algclosed characteristic zero BFJ plus DVR/empty cases. No catalogue or external acceptance is performed.
